\documentclass[../main.tex]{subfiles}
\begin{document}

\chapter{Glossary}
\label{ch:glossary}

Each entry gives a one-line definition and a pointer to the section where the term is treated in full. For page-level lookup of any bold term in the text, use the index at the very back of the guide.

\begin{description}[leftmargin=0pt, style=unboxed]

\glentry{3-2-1 rule}{A backup standard of three copies of the data, on two different storage media, with one copy in a separate physical location (see \autoref{sec:dmp}).}

\glentry{AIAA style (numbered)}{A citation style that marks sources with bracketed numbers in order of appearance and lists them numerically, keeping the text clean but uninformative in itself (see \autoref{sec:referencing}).}

\glentry{APA style (author-date)}{A citation style that names author and year in the text and lists sources alphabetically, letting a reader recognise key works at the cost of longer in-text references (see \autoref{sec:referencing}).}

\glentry{Attribute}{A value that a variable can take, which may be quantitative or qualitative and must be exhaustive and mutually exclusive (see \autoref{sec:language-research}).}

\glentry{Causal study}{A study designed to determine whether one or more variables cause or affect one or more outcome variables (see \autoref{sec:language-research}).}

\glentry{Complex probability sampling}{Probability sampling that improves efficiency by dropping the equal-chance requirement while keeping selection probabilities known, as in systematic, cluster, and multi-stage designs (see \autoref{sec:sample-types}).}

\glentry{Conceptual research design}{The strategic area that models the content of the research and fixes what, why, and how much is studied (see \autoref{sec:conceptual-technical}).}

\glentry{Convenience sampling}{Non-probability sampling in which units are chosen because they are available, which must be stated and cannot be assumed representative (see \autoref{sec:sample-types}).}

\glentry{Correlation versus causation}{The principle that an observed association between variables does not establish that one drives the other, since a third variable may drive both (see \autoref{sec:language-research}).}

\glentry{Critical path}{The longest chain of dependent tasks in a schedule, which sets the shortest possible duration of the whole project (see \autoref{sec:project-mgmt}).}

\glentry{Data management plan (DMP)}{A six-section plan, designed before collection, that describes, documents, stores, clears, shares, and assigns responsibility for a project's data (see \autoref{sec:dmp}).}

\glentry{Data quality}{The set of checks run before analysis for typing errors, outliers, coding errors, and agreement with expectations (see \autoref{sec:data-sources}).}

\glentry{Deduction}{Reasoning that moves from the more general to the more specific, forming the narrowing half of the hourglass (see \autoref{sec:hourglass}).}

\glentry{Descriptive study}{A study designed to document the values of one or more unknown variables (see \autoref{sec:language-research}).}

\glentry{Engineering for One Planet (EOP)}{A framework that judges research on its value, its reproducibility, and its resource management across environmental, social, cultural, and economic implications (see \autoref{sec:eop}).}

\glentry{FAIR principles}{The guiding principles that data should be findable, accessible, interoperable, and reusable, with reuse as the ultimate objective (see \autoref{sec:dmp}).}

\glentry{Feasibility}{The trade-off between rigour and practicalities, covering duration, ethics, cooperation, cost, and access, that gates whether an idea becomes a project (see \autoref{sec:research-ideas}).}

\glentry{Gantt chart}{A time-based chart of a project's tasks that schedules the work and exposes the critical path (see \autoref{sec:project-mgmt}).}

\glentry{Hourglass model}{The shape of a single study, narrowing from broad questions through operationalisation and measurement, then widening back to generalised conclusions (see \autoref{sec:hourglass}).}

\glentry{Human Research Ethics Committee (HREC)}{The TU Delft board whose permission is required for any research involving human test subjects (see \autoref{sec:ethics}).}

\glentry{Hypothesis}{A specific statement of prediction describing in concrete terms what a study expects to happen, whose form fixes the study type (see \autoref{sec:language-research}).}

\glentry{Induction}{Reasoning that moves from specific observations to broader generalisations, the fallible widening half of the hourglass (see \autoref{sec:hourglass}).}

\glentry{Literature review}{A survey that identifies related research and sets the project within a conceptual and theoretical context, locating the gap to be filled (see \autoref{sec:research-ideas}).}

\glentry{Operationalisation}{The conceptual step that turns abstract concepts into exact definitions, stated assumptions, named variables, and a measurement plan (see \autoref{sec:operationalization}).}

\glentry{Personal data}{Any data traceable directly or indirectly to a natural person, which falls under the General Data Protection Regulation (see \autoref{sec:data-sources}).}

\glentry{Population}{The entire collection of things in a field, the whole problem restated as a collection and too large to study directly (see \autoref{sec:sampling}).}

\glentry{Probability sampling}{Sampling in which every element has a known probability of selection, the precondition for statistical generalisation (see \autoref{sec:sample-types}).}

\glentry{Protected data}{Knowledge in fields such as nuclear weapons and missile technology, subject to a knowledge embargo requiring an exemption (see \autoref{sec:data-sources}).}

\glentry{Purposive sampling}{Non-probability sampling in which units are chosen because they hold desired information, which must be stated and cannot be assumed representative (see \autoref{sec:sample-types}).}

\glentry{Qualitative data}{Data collected by observation of an environment, such as text, photographs, and people (see \autoref{sec:data-sources}).}

\glentry{Quantitative data}{The numerical result of measured or simulated variables, such as wind-tunnel data, finite-element results, and surveys (see \autoref{sec:data-sources}).}

\glentry{Relational study}{A study investigating the relationship between two or more variables (see \autoref{sec:language-research}).}

\glentry{Reproducibility}{The property that work can be repeated by others, signalling rigour and enabling reuse with fewer resources (see \autoref{sec:eop}).}

\glentry{Research framework}{A schematic representation of the steps needed to answer the research questions, naming the required knowledge and methods (see \autoref{sec:objectives-rqs}).}

\glentry{Research logbook}{A detailed record of design choices, assumptions, procedures, incidents, file histories, and people consulted that makes a project traceable (see \autoref{sec:data-sources}).}

\glentry{Research objective}{A statement of the contribution a project makes to a problem and the way that contribution is delivered, judged informative, useful, feasible, and clear (see \autoref{sec:objectives-rqs}).}

\glentry{Research question}{An open-ended question, one of at least two, that yields information needed for the objective and enables the technical design (see \autoref{sec:objectives-rqs}).}

\glentry{Research strategy}{The technical-design phase that defines exactly how the research will be done, judged on suitability, availability, duration, and ethics (see \autoref{sec:strategy-materials}).}

\glentry{Resource management}{Awareness and control of the people, equipment, facilities, and digital resources a project consumes, split into digital resources and methodologies (see \autoref{sec:eop}).}

\glentry{Sample}{The chosen subset of a population that is actually studied (see \autoref{sec:sampling}).}

\glentry{Sampling frame}{The concrete means of listing and accessing the study population from which the sample is drawn (see \autoref{sec:sampling}).}

\glentry{Simple random sampling}{Drawing so that every element is known and every possible sample has an equal probability of selection (see \autoref{sec:sample-types}).}

\glentry{Standard deviation}{The spread of scores around their average in a single sample, the square root of the variance, in the same units as the measure (see \autoref{sec:sample-types}).}

\glentry{Standard error}{The precision of a statistical sample, calculated from the standard deviation, falling as the sample grows and the deviation shrinks (see \autoref{sec:sample-types}).}

\glentry{Study population}{The part of the theoretical population that can actually be accessed (see \autoref{sec:sampling}).}

\glentry{Technical report}{A document that conveys specific information about a technical subject to a specific audience for a specific and practical purpose (see \autoref{sec:reports}).}

\glentry{Technical research design}{The strategic area that plans the execution of the research and fixes how, where, and when through strategy, materials, and planning (see \autoref{sec:conceptual-technical}).}

\glentry{Theoretical population}{The ideal group to which a study wants to generalise, the starting point of the sampling funnel (see \autoref{sec:sampling}).}

\glentry{Validation}{Proving that the outcomes of the research are true and based on strong scientific evidence, an outcome-facing and often external check (see \autoref{sec:vv}).}

\glentry{Value of research}{The EOP axis that requires forecasting the near-term and long-term benefit and potential harm of the work (see \autoref{sec:eop}).}

\glentry{Variable}{The smallest block in the language of research, whose possible values are its attributes (see \autoref{sec:language-research}).}

\glentry{Verification}{Proving that the research method has been used correctly and is suitable for the topic, a process-facing and often internal check (see \autoref{sec:vv}).}

\glentry{Work Breakdown Structure (WBS)}{A decomposition of a project into work packages and tasks, taken to a level of detail that allows each task to be estimated (see \autoref{sec:project-mgmt}).}

\end{description}

\end{document}
